\documentclass[10pt,a4paper]{amsart}
\usepackage[margin=19mm]{geometry}
\usepackage{amsmath,amssymb,mathtools}
\usepackage[scr=rsfs,cal=euler]{mathalfa}
\usepackage{xfrac}
\usepackage[unicode,hidelinks,bookmarks=false]{hyperref}
\newcommand{\ie}{i.e.\ }
\newcommand{\cf}{cf.\ }
\newcommand{\resp}{respectively\ }
\newcommand{\Subsection}{\subsection}
\providecommand{\nobreakdash}{\nobreak\hbox{-}}
\setlength{\emergencystretch}{2em}
\multlinegap0pt
\usepackage{bm}
\usepackage{bbm}
\usepackage{dsfont}
\usepackage{xspace}
\usepackage{cancel}
\usepackage{mathtools}
\usepackage{enumitem}
\usepackage{amsthm}
\makeatletter
\@ifundefined{theorem}{\newtheorem{theorem}{Theorem}}{}
\@ifundefined{proposition}{\newtheorem{proposition}[theorem]{Proposition}}{}
\makeatother
\title[A Unified Substitution Method for Integration]{A Unified Substitution Method for Integration}
\author{Emmanuel Antonio Jos\'e Garc\'ia}
\date{Source: arXiv:2505.03754v5 (CC BY 4.0)}
\begin{document}
\maketitle

\noindent\emph{Source and licence.} A Unified Substitution Method for Integration. Version arXiv:2505.03754v5 (CC BY 4.0), distributed under the Creative Commons Attribution 4.0 International licence, as recorded in the arXiv licence metadata for this exact version and shown on the abstract page https://arxiv.org/abs/2505.03754v5. Full rights notice: licenses/arxiv-2505.03754-CC-BY-4.0.txt.\par\smallskip

\noindent\textit{Editorial scope.} Counted unit: the Proposition whose source label is \texttt{prop:laurent} in the article (source0.tex lines 781--827, source label \texttt{prop:laurent}), delivered together with the article's own complete original proof of it (source0.tex lines 829--908, one \texttt{proof} environment of 80 lines). No mathematics is written, repaired or re-derived: statement and proof are the article's own text line for line. Required context retained uncounted: none. Every definition, display and label of those lines is retained unchanged. Omitted from this package and not redistributed: 1--780, 828--828, 909--1876. Nothing inside a retained excerpt is omitted or altered: every retained body is delivered exactly as the article's lines are written, so no omission receipt of any kind accompanies this package. Citations: the retained excerpts carry no citation command at all, so no bibliography entry of the article is transcribed and no reference is redistributed. Reference adaptation: none was needed and none was made; every reference inside the delivered unit resolves inside it, so no unresolved reference or citation is produced. Printed slips kept literal: no printed word, symbol, hypothesis or number of the counted statement or of its proof was altered. Layout and class: the article's own class is \texttt{article} with packages including geometry, fontenc, inputenc, lmodern, microtype, amsmath, amssymb, amsthm, mathtools, booktabs, enumitem, hyperref, xurl, array; the wrapper is the lane's pinned \texttt{amsart} editorial wrapper at 10pt on A4, which restates the theorem-like environments the retained text needs and adds the article's own macro definitions that the retained text needs (none), with extra wrapper packages (bm, bbm, dsfont, xspace, cancel, mathtools, enumitem). The delivered numbering is the unit's own. Assets: no retained span contains a figure environment, a graphics-inclusion command, a PDF-page inclusion, a table or a TikZ picture; no image file is copied into this package, the package declares no asset dependency and no image file is redistributed with it. Licence: version arXiv:2505.03754v5 of this article is distributed under the Creative Commons Attribution 4.0 International licence, as recorded in the arXiv licence metadata for this exact version and shown on the abstract page https://arxiv.org/abs/2505.03754v5. A full rights notice is delivered at licenses/arxiv-2505.03754-CC-BY-4.0.txt. Characters: every retained excerpt is pure ASCII, so the delivered engine, XeLaTeX with its default Latin Modern text font, reports no missing character.
\begin{proposition}[Laurent-polynomial families]
\label{prop:laurent}
Let $X=x+b$, let $P$ be a polynomial, and restrict all expressions to an open component on which the corresponding transform and principal-root representation hold and every displayed denominator is nonzero.
\begin{enumerate}[label=\arabic*.]
\item Under Transform~2, for every polynomial $P$,
\[
P(x)\sqrt{X^2-a^2}\,dx
\]
becomes a finite Laurent polynomial in $t$.
\item Under Transform~2, for every polynomial $P$,
\[
P(x)\sqrt{\frac{X-a}{X+a}}\,dx
\]
becomes a finite Laurent polynomial in $t$.
\item Under Transform~3, for every polynomial $P$, both
\[
P(x)\sqrt{X^2+a^2}\,dx,
\qquad
\frac{P(x)}{\sqrt{X^2+a^2}}\,dx
\]
become finite Laurent polynomials in $s$. Moreover, for every integer $n\ge1$, each of
\[
\frac{dx}{\left(X+\sqrt{X^2+a^2}\right)^n},
\qquad
\frac{dx}{\left(\sqrt{X^2+a^2}-X\right)^n}
\]
becomes a finite Laurent polynomial in $s$.
\item Under Transform~5, for every integer $n\ge1$,
\[
\frac{dx}{X^n\sqrt{a^2-X^2}}
\]
becomes a finite Laurent polynomial in $r$.
\item Under Transform~5, for every integer $n\ge3$,
\[
\frac{\sqrt{a^2-X^2}}{X^n}\,dx
\]
becomes a finite Laurent polynomial in $r$.
\item Under Transform~5, for every integer $n\ge2$, each of
\[
\frac{1}{X^n}\sqrt{\frac{a+X}{a-X}}\,dx,
\qquad
\frac{1}{X^n}\sqrt{\frac{a-X}{a+X}}\,dx
\]
becomes a finite Laurent polynomial in $r$.
\end{enumerate}
When $b=0$, one has $X=x$, so the last three families are precisely the corresponding families with reciprocal powers of $x$.
\end{proposition}

\begin{proof}
For Transform~2,
\[
x=\frac{a}{2}\left(t+t^{-1}\right)-b,
\]
so every polynomial $P(x)$ is a Laurent polynomial in $t$. The difference radical gives
\[
\sqrt{X^2-a^2}\,dx
=-\varepsilon\frac{a^2}{4}t^{-3}(t^2-1)^2\,dt,
\]
which is Laurent because $\varepsilon$ is constant on each component. The radical quotient cancels the factor $1+t$ in the Jacobian:
\begin{align*}
\sqrt{\frac{X-a}{X+a}}\,dx
&=\frac{1-t}{1+t}\,a\frac{t^2-1}{2t^2}\,dt\\
&=-\frac{a}{2}\left(1-t^{-1}\right)^2\,dt.
\end{align*}
Multiplication by $P(x)$ therefore gives a Laurent polynomial in both cases.

For Transform~3, write
\[
H:=\sqrt{X^2+a^2}.
\]
Then
\[
x=\frac{a}{2}\left(s-s^{-1}\right)-b,
\qquad
H=\frac{a}{2}\left(s+s^{-1}\right),
\qquad
dx=\frac{a}{2}\left(1+s^{-2}\right)\,ds.
\]
Thus every polynomial $P(x)$ is a Laurent polynomial in $s$, and
\[
H\,dx=\frac{a^2}{4}\left(s+2s^{-1}+s^{-3}\right)\,ds,
\qquad
\frac{dx}{H}=s^{-1}\,ds.
\]
Consequently, both $P(x)H\,dx$ and $P(x)H^{-1}\,dx$ are finite Laurent polynomials in $s$. In addition,
\[
X+H=as,
\qquad
H-X=as^{-1},
\]
and hence, for $n\ge1$,
\[
\frac{dx}{(X+H)^n}
=\frac{a^{1-n}}{2}\left(s^{-n}+s^{-n-2}\right)\,ds,
\]
\[
\frac{dx}{(H-X)^n}
=\frac{a^{1-n}}{2}\left(s^n+s^{n-2}\right)\,ds.
\]
Each expression is a finite Laurent polynomial.

For Transform~5,
\[
X=\frac{2ar}{1+r^2},
\qquad
\sqrt{a^2-X^2}=a\frac{1-r^2}{1+r^2},
\qquad
dx=2a\frac{1-r^2}{(1+r^2)^2}\,dr.
\]
Direct cancellation gives
\[
\frac{dx}{X^n\sqrt{a^2-X^2}}
=2^{1-n}a^{-n}r^{-n}(1+r^2)^{n-1}\,dr,
\qquad n\ge1,
\]
\[
\frac{\sqrt{a^2-X^2}}{X^n}\,dx
=2^{1-n}a^{2-n}r^{-n}(1-r^2)^2(1+r^2)^{n-3}\,dr,
\qquad n\ge3,
\]
and
\[
\frac{1}{X^n}\sqrt{\frac{a\pm X}{a\mp X}}\,dx
=2^{1-n}a^{1-n}r^{-n}(1\pm r)^2(1+r^2)^{n-2}\,dr,
\qquad n\ge2.
\]
Each right-hand side is a finite Laurent polynomial. The stated restrictions on $n$ are exactly those that prevent a residual negative power of $1+r^2$ in these formulas.
\end{proof}

\end{document}
